\chapter{Minimalne przypomnienie geometrii różniczkowej}
\label{ch:dodatek-geometria}

W części poświęconej kohomologii de Rhama, dualności Poincarégo i klasom
charakterystycznym potrzebujemy pojęć geometrycznych. Ten dodatek
nie zastępuje dowodów z książki o geometrii różniczkowej, ale wyjaśnia
terminy oraz precyzuje przyjmowane twierdzenia.

\section{Rozmaitości gładkie, przestrzeń styczna i formy}
\begin{definicja}
Gładka rozmaitość wymiaru $n$ jest rozmaitością topologiczną Hausdorffa,
spełniającą drugi aksjomat przeliczalności, z atlasem map do otwartych
podzbiorów $\R^n$, w którym zmiany współrzędnych są gładkie.
Przestrzeń styczna $T_pM$ w punkcie $p$ jest przestrzenią derywacji
$C^\infty(M)\to\R$ w punkcie $p$; przestrzeń kostyczna to dualna
$T_p^*M$. Formy różniczkowe stopnia $k$ są gładkimi przekrojami
wiązki $\bigwedge^kT^*M$.
\end{definicja}
\begin{stwierdzenie}
Jeżeli $V$ jest $n$-wymiarową przestrzenią wektorową z bazą
$e_1,\dots,e_n$, a $e^1,\dots,e^n$ jest bazą dualną,
to dla każdego $0\leq k\leq n$ formy
$e^{i_1}\wedge\cdots\wedge e^{i_k}$,
$1\leq i_1<\cdots<i_k\leq n$, tworzą bazę $\bigwedge^k V^*$.
\end{stwierdzenie}
\begin{proof}
Z naprzemienności każdy prosty iloczyn jest równy zeru lub
odpowiedniemu iloczynowi o ściśle rosnących indeksach z właściwym znakiem,
więc wskazane elementy rozpinają $\bigwedge^kV^*$.
Ich niezależność wynika przez wartościowanie na zestawach
$(e_{j_1},\dots,e_{j_k})$: dany iloczyn ma wartość $1$ dla
odpowiadającego zestawu i $0$ dla pozostałych.
\end{proof}

\section{Regularne podrozmaitości, transwersalność i otoczenia tubularne}
\begin{definicja}
Gładkie odwzorowanie $f:M\to N$ jest submersją w $p$, gdy
$\mathrm df_p:T_pM\to T_{f(p)}N$ jest surjekcją.
Podrozmaitości $A,B\subseteq M$ są transwersalne w $p\in A\cap B$,
jeśli $T_pA+T_pB=T_pM$.
Jeśli $S\subset M$ jest podrozmaitością osadzoną,
\emph{wiązką normalną} $\nu(S,M)$ nazywamy $TM|_S/TS$.
\end{definicja}
\begin{twierdzenie}[Zestaw twierdzeń wykorzystywanych dalej]
Jeżeli $f:M\to N$ jest submersją, to w odpowiednich lokalnych
współrzędnych ma postać projekcji $(x,y)\mapsto x$.
Jeżeli dwie gładkie podrozmaitości $A,B$ przecinają się transwersalnie,
to $A\cap B$ jest podrozmaitością o wymiarze $\dim A+\dim B-\dim M$.
Dla osadzonej podrozmaitości $S\subset M$ istnieje dyfeomorfizm
otoczenia przekroju zerowego w $\nu(S,M)$ na otoczenie $S$ w $M$.
Dla gładkich rozmaitości spełnione są standardowe twierdzenia
Sarda i transwersalności.
\end{twierdzenie}
\begin{uwaga}
Twierdzenia te są tu przyjmowane z wcześniejszej książki.
W szczególności użycie transwersalności przy liczeniu przecięć
wymaga zwartości lub odpowiednich warunków właściwości dla
zapewnienia skończoności zbioru punktów przecięcia.
\end{uwaga}

\begin{definicja}
\emph{Uramowaniem} rzeczywistej wiązki wektorowej rzędu $r$
nazywamy wybór $r$ ciągłych przekrojów, które w każdym włóknie
tworzą bazę. Istnienie uramowania jest równoważne trywialności wiązki.
\end{definicja}

\section{Gdzie kontynuować dalej?}
Po opanowaniu teorii zawartej w skrypcie naturalnymi dalszymi etapami są
silne homologie, dalsza teoria ciągów spektralnych
(w szczególności ciąg Adamsa), operacje wtórne,
obliczenia pierścieni kobordyzmu oraz bardziej zaawansowana teoria widm.
Jeszcze dalej leżą torsja Whiteheada, teoria chirurgii
i topologia czterowymiarowa; wymagają znacznie większego przygotowania.

